\documentclass[11pt,a4paper]{article}
\usepackage[T2A]{fontenc}
\usepackage[utf8]{inputenc}
\usepackage[russian]{babel}
\usepackage{amsmath,amssymb}
\usepackage[a4paper,margin=20mm,headheight=14pt,footskip=14mm]{geometry}
\usepackage{xcolor}
\usepackage{enumitem}
\usepackage{fancyhdr}
\usepackage{microtype}
\usepackage{needspace}
\usepackage{array}

\definecolor{ink}{HTML}{223142}
\definecolor{accent}{HTML}{176B83}
\definecolor{soft}{HTML}{F1F7F9}
\definecolor{line}{HTML}{CBD8DE}
\setlength{\parindent}{0pt}
\setlength{\parskip}{5pt}
\setlist[enumerate]{leftmargin=*,itemsep=3pt,topsep=4pt}
\renewcommand{\baselinestretch}{1.09}
\pagestyle{fancy}
\fancyhf{}
\renewcommand{\headrulewidth}{0pt}
\fancyfoot[C]{\footnotesize\color{ink} Лёвин Артём Александрович эксклюзивно для Ксении Васильченко}
\fancyfoot[R]{\footnotesize\color{ink}\thepage}
\newcommand{\problemtitle}[2]{%
  \Needspace{7\baselineskip}%
  \vspace{8pt}{\large\bfseries\color{accent} Задание №~#1. #2}\par\vspace{2pt}
  {\color{line}\hrule height 0.7pt}\vspace{6pt}}
\newcommand{\practice}{\vspace{4pt}\textbf{\color{accent} Задача на закрепление.}\quad}
\begin{document}
\begin{center}
  {\Large\bfseries\color{ink} Проверка домашней работы}\\[4pt]
  {\normalsize Подробный разбор ошибок}\\[7pt]
  {\bfseries Ксения Васильченко}\quad\textbar\quad 22 октября 2026 г.
\end{center}
\vspace{3pt}
\noindent\fcolorbox{line}{soft}{\parbox{\dimexpr\linewidth-2\fboxsep-2\fboxrule}{%
\textbf{Результат проверки.} Правильно выполнены задания №~1--5, 8 и 9. Ошибки допущены в заданиях №~6 и 7. Задание №~10 содержит верную схему, но вычисления и ответ отсутствуют. Ниже разобраны только задания, требующие доработки.}}

\problemtitle{6}{Действия с обыкновенными дробями}
\textbf{Условие:}
\[
\left(\frac49+\frac5{12}\right)\cdot\frac{18}{31}.
\]
\textbf{Где возникла ошибка.} В тетради после первого знака равенства записана только сумма дробей, хотя исходное выражение содержит ещё умножение на \(\frac{18}{31}\). Далее преобразования с числителями и знаменателями приводят к неверному результату. Сокращать числители и знаменатели отдельных слагаемых между собой при сложении нельзя: сокращение применяется к \emph{множителям} числителя и знаменателя одной дроби.

\textbf{Правильное решение по шагам.}
\begin{enumerate}
\item \textbf{Сначала вычисляем сумму в скобках.} Наименьший общий знаменатель дробей со знаменателями 9 и 12 равен 36. Поэтому
\[
\frac49=\frac{4\cdot4}{9\cdot4}=\frac{16}{36},\qquad
\frac5{12}=\frac{5\cdot3}{12\cdot3}=\frac{15}{36}.
\]
\item \textbf{Складываем числители, общий знаменатель сохраняем:}
\[
\frac{16}{36}+\frac{15}{36}=\frac{31}{36}.
\]
\item \textbf{Теперь умножаем на дробь, стоящую после скобок:}
\[
\frac{31}{36}\cdot\frac{18}{31}
=\frac{31\cdot18}{36\cdot31}
=\frac{18}{36}=\frac12.
\]
\end{enumerate}
\textbf{Ответ:} \(\boxed{\frac12}\).

\textbf{Что запомнить.} При наличии скобок сначала выполняем действие внутри них. Для сложения дробей нужен общий знаменатель. Сокращение при умножении допустимо, когда одинаковый множитель находится в числителе и знаменателе.

\practice Вычислите:
\[
\left(\frac38+\frac5{12}\right)\cdot\frac{24}{19}.
\]

\problemtitle{7}{Смешанные числа и сокращение дроби}
\textbf{Условие:}
\[
\left(2\frac34-1\frac14\right)\cdot\frac37.
\]
\textbf{Где возникла ошибка.} В тетради верно выполнены перевод смешанных чисел в неправильные дроби, вычитание и умножение. Получена дробь \(\frac{18}{28}\). Ошибка появилась в последнем действии: эта дробь была сокращена до \(\frac37\). При сокращении необходимо делить \emph{числитель и знаменатель на одно и то же число}. Здесь такой переход неверен.

\textbf{Правильное решение по шагам.}
\begin{enumerate}
\item \textbf{Переводим смешанные числа в неправильные дроби:}
\[
2\frac34=\frac{2\cdot4+3}{4}=\frac{11}{4},\qquad
1\frac14=\frac{1\cdot4+1}{4}=\frac54.
\]
\item \textbf{Вычитаем дроби с одинаковыми знаменателями:}
\[
\frac{11}{4}-\frac54=\frac{11-5}{4}=\frac64.
\]
\item \textbf{Умножаем полученную дробь на \(\frac37\):}
\[
\frac64\cdot\frac37=\frac{6\cdot3}{4\cdot7}=\frac{18}{28}.
\]
\item \textbf{Сокращаем числитель и знаменатель на 2:}
\[
\frac{18}{28}=\frac{18:2}{28:2}=\frac9{14}.
\]
\end{enumerate}
\textbf{Ответ:} \(\boxed{\frac9{14}}\).

\textbf{Что запомнить.} Чтобы сократить дробь, находим общий делитель числителя и знаменателя. Делим оба числа на \emph{один и тот же} делитель. Полезно проверить: числитель результата должен оставаться меньше знаменателя, поскольку исходное произведение меньше единицы.

\practice Вычислите:
\[
\left(3\frac14-1\frac34\right)\cdot\frac49.
\]

\clearpage
\problemtitle{10}{Текстовая задача о прыжках}
\textbf{Условие.} Длина прыжка жабы равна \(1\frac15\) см. Прыжок лягушки в \(2\frac13\) раза длиннее. Оба животных одновременно прыгают из одного места в одном направлении и делают по 15 прыжков. Определите расстояние между ними.

\textbf{Что требует доработки.} Схема в тетради отражает условие верно: оба начинают из одной точки, прыжок лягушки длиннее, а расстояние между конечными положениями обозначено знаком вопроса. Числовое решение и ответ отсутствуют. Для полного решения нужно вычислить длину прыжка лягушки и сравнить расстояния после 15 прыжков.

\textbf{Правильное решение по шагам.}
\begin{enumerate}
\item \textbf{Переводим смешанные числа в неправильные дроби:}
\[
1\frac15=\frac65,\qquad 2\frac13=\frac73.
\]
\item \textbf{Находим длину одного прыжка лягушки.} Её прыжок длиннее в \(2\frac13\) раза, поэтому длину прыжка жабы умножаем на это число:
\[
\frac65\cdot\frac73=\frac{42}{15}=\frac{14}{5}=2\frac45\text{ см}.
\]
\item \textbf{Находим расстояние, пройденное жабой за 15 прыжков:}
\[
\frac65\cdot15=18\text{ см}.
\]
\item \textbf{Находим расстояние, пройденное лягушкой за 15 прыжков:}
\[
\frac{14}{5}\cdot15=42\text{ см}.
\]
\item \textbf{Вычитаем пройденные расстояния,} так как животные начали движение вместе и прыгали в одном направлении:
\[
42-18=24\text{ см}.
\]
\end{enumerate}
\textbf{Ответ:} \(\boxed{24\text{ см}}\).

\textbf{Как проверить себя другим способом.} За один прыжок лягушка удаляется от жабы на
\[
\frac{14}{5}-\frac65=\frac85\text{ см}.
\]
За 15 прыжков расстояние между ними станет \(\frac85\cdot15=24\) см. Ответ совпал.

\practice Жаба прыгает на \(1\frac14\) см, а лягушка --- в \(2\frac15\) раза дальше за один прыжок. Обе начали прыгать из одной точки в одном направлении. Определите расстояние между ними после 12 прыжков каждой.

\vfill
\begin{center}\small\color{ink} \textit{После разбора решите три задания на закрепление самостоятельно.}\end{center}
\end{document}
